\documentstyle[%


righttag%


,11pt%


,amssymb%


,russian%               remove in Eng.


,izvru%                 Set izven% in Eng.


]{amsart}





% seek the \text constructions


\newcommand{\tts}[1]{}


\numberwithin{equation}{section}


\renewcommand{\baselinestretch}{1.03}





%\hoffset=-15mm%                  For Eng.


\setcounter{page}{21}


\god{2002}


\nomer{1~(476)} %                        Remove in Eng.


%\nomer{Vol.\,46, No.\,1}%               Set in Eng.


%\str{\pageref{begin}--\pageref{end}}%  Set in Eng.


\udk{517.988.8}





\author{\it С.Е.\,Железовский}


% NB:   ^^ remove in Eng.


\title{Оценки скорости сходимости


проекционно-разностного метода


для гиперболических уравнений}





\date{}


%\pagestyle{myheadings}%         Set in Eng.


\pagestyle{plain} %              Remove in Eng.


\begin{document}


%\thispagestyle{plain}%          Set in Eng.


\maketitle


\label{begin}








\section*{Введение}





В данной статье исследуется сходимость схемы решения абстрактного


гиперболического уравнения в гильбертовом пространстве с переменными


операторными коэффициентами, являющейся комбинацией схемы Гал\"еркина по


пространству и трехслойной разностной схемы с весами по


времени. Какие-либо специальные условия на конечномерные подпространства,


в которых должны принимать значения приближенные решения, отсутствуют.


Устанавливаются асимптотические энергетические оценки скорости сходимости


рассматриваемой вычислительной схемы. В конце статьи показано, что полученные


результаты применимы к исследованию сходимости схем проекционно-разностного


метода решения начально-краевых задач для гиперболических уравнений второго


порядка с дискретизацией по пространству методом конечных элементов.





Основные результаты работы содержатся в теоремах 4--6. Теорема 4 относится


к случаю естественной гладкости точного решения исходного уравнения, а теоремы


5, 6 --- к случаям повышенной гладкости точного решения. В двух последних


теоремах даны оценки, неулучшаемые по порядку. Близкие результаты для


гиперболических уравнений второго порядка с коэффициентами, не зависящими


от времени, и симметричных схем были получены в [1]. Для полудискретного


метода Гал\"еркина


оценки, аналогичные установленным здесь, были получены в [2]. Для абстрактных


параболических уравнений в случае чисто неявной разностной схемы по времени


результаты, подобные результатам данной работы, устанавливались в [3], [4].





В статье используются стандартные обозначения вида


$\cal L(X;Y)$, $C^k([0,T];X)$, $L^p(0,T;X)$, где $X$ и $Y$ ---


линейные нормированные пространства, $[0,T]$ --- отрезок оси ${\bold R}$


(напр., [5], с.\,579), а также следующие обозначения: $\cal B(X)$ ---


пространство билинейных форм $\varphi:X\times X\to{\bold R}$, для которых


$\|\varphi\|_{\cal B(X)}=


\sup\limits_{\|x_1\|_X=\|x_2\|_X=1}|\varphi(x_1,x_2)|


<\infty$; $W^{k,p}(0,T;X)$, $k\in{\bold N}$, $p=1$ или $p=\infty$,


--- пространство функций из $C^{k-1}([0,T];X)$, производные $k$-го порядка


которых, понимаемые в смысле распределений ([5], с.\,20), принадлежат


$L^p(0,T;X)$. Кроме того, используются обозначения вида $\wh y$, $\check y$,


$y_t$, $y_{\ov t}$, $y_{\overset{\,\circ}{t}}$, $y_{\ov tt}$, принятые


в теории разностных схем ([6], cс.\,613, 614; [7], с.\,10).





\section{Исходная задача}





Пусть $H$ --- вещественное сепарабельное гильбертово пространство со скалярным


произведением $(\cdot,\cdot\cdot)_H$ и нормой $\|\cdot\|_H$, $T>0$ --- заданное


число. Будем рассматривать задачу Коши для абстрактного гиперболического


уравнения


\begin{gather}


u''(t)+A(t)u(t)+B(t)u(t)=f(t),\quad t\in[0,T],\\


u(0)=u'(0)=0,


\end{gather}


где $u:[0,T]\to H$ --- искомая функция, $A(t)$, $B(t)$, $t\in[0,T]$, ---


заданные операторы, $f:[0,T]\to H$ --- заданная функция. Все операторы


$A(t)$, $t\in[0,T]$, --- неограниченные, самосопряженные и положительно


определенные в $H$ и имеют общую область определения $D\subset H$; оператор


$(A(0))^{-1}$, обратный к $A(0)$, компактен в $H$. В силу теоремы Гайнца


([8], с.\,177) все операторы $(A(t))^{1/2}$, $t\in[0,T]$, также имеют общую


область определения, которую обозначим через $V$. Введем на $D$ норму


$\|\cdot\|_D=\|A(0)\cdot\|_H$, тем самым превратим $D$ в банахово пространство.


Примем, что для любых $v\in D$ и $t\in[0,T]$


\begin{gather}


%%%%%% return to align


\|v\|_D\le a_1\|A(t)v\|_H,\\


\|A(t)v\|_H\le a_2\|v\|_D,


\end{gather}


где $a_1$, $a_2$ --- константы, не зависящие от $v$ и $t$. Обозначим


энергетическое скалярное произведение и энергетическую норму, порожденные


оператором $A(t)$, соответственно через $a(t,\cdot,\cdot\cdot)$ и


$\|\cdot\|_{(t)}$, т.\,е. $a(t,v,w)=((A(t))^{1/2}v,(A(t))^{1/2}w)_H$,


$\|v\|_{(t)}=\|(A(t))^{1/2}v\|_H$,


$t\in[0,T]$, $v,w\in V$. Снабдив $V$ скалярным произведением


$(\cdot,\cdot\cdot)_V=a(0,\cdot,\cdot\cdot)$ и соответственно нормой


$\|\cdot\|_V=\|\cdot\|_{(0)}$,


превратим это множество в сепарабельное гильбертово пространство.


Будем считать, что


функция ``$t\to a(t,\cdot,\cdot\cdot)$'' принадлежит пространству


$W^{2,1}(0,T;\cal B(V))$, функция ``$t\to B(t)$'' --- пространству


$W^{1,1}(0,T;\cal L(V;H))$, a


$f\in W^{1,1}(0,T;H)$.





Все сформулированные здесь условия всюду в дальнейшем считаются выполненными


без специальных оговорок. Все дополнительные условия, налагаемые ниже на


операторы $A(t)$, $B(t)$ и функцию $f$, считаются выполненными только тогда,


когда об этом сказано явно.





Отметим для дальнейшего некоторые следствия принятых условий. Из (1.3), (1.4)


по теореме Гайнца следует, что для любых $v\in D$ и $t\in[0,T]$


\begin{gather}


%%%%% return to align


\|v\|_V\le a_1^{1/2}\|v\|_{(t)},\\


\|v\|_{(t)}\le a_2^{1/2}\|v\|_V.


\end{gather}


Так как ``$t\to a(t,\cdot,\cdot\cdot)\text{''}\in W^{2,1}(0,T;\cal B(V))$,


то заведомо ``$t\to a(t,\cdot,\cdot\cdot)\text{''}\in C^1([0,T];\cal B(V))$,


поэтому существует такая константа $a_3$, что для любых $t\in[0,T]$,


$v,w\in V$


\begin{equation}


|a'(t,v,w)|\le a_3\|v\|_V\|w\|_V,


\end{equation}


где ``$t\to a'(t,\cdot,\cdot\cdot)$'' --- производная первого порядка функции


``$t\to a(t,\cdot,\cdot\cdot)$''. Так как


``$t\to B(t)\text{''}\in W^{1,1}(0,T;\cal L(V;H))$, то заведомо


``$t\to B(t)\text{''}\in C^0([0,T];\cal L(V;H))$,


поэтому существует такая константа $b_0$, что для любых


$t\in[0,T]$, $v\in V$


\begin{equation}


\|B(t)v\|_H\le b_0\|v\|_V.


\end{equation}





\begin{thm}


Задача $(1.1)$, $(1.2)$ имеет единственное


решение


$u\in W^{2,\infty}(0,T;H)\cap\linebreak W^{1,\infty}(0,T;V)\cap


L^{\infty}(0,T;D)$,


удовлетворяющее уравнению $(1.1)$ при почти всех $t\in[0,T]$


и начальным условиям $(1.2)$ непосредственно.


\end{thm}





Доказательство существования решения задачи (1.1), (1.2) в указанном смысле


проводится по известной схеме с использованием полудискретных аппроксимаций


Гал\"еркина ([5], cс.\,22--27, 30--32; [9], с.213--220). Доказательство


единственности решения аналогично, например, доказательству теоремы 2.1


в ([9], гл.\,IV, с.\,206).





Ниже в теоремах 5, 6 предполагается, что решение $u$ задачи (1.1), (1.2)


обладает определенной повышенной гладкостью. Требуемая в этих теоремах


(фактически несколько более высокая) гладкость решения задачи (1.1), (1.2)


имеет место при надлежащем усилении условий на операторы $A(t)$, $B(t)$


и функцию $f$. Сформулируем здесь соответствующие результаты.





\begin{thm}


Пусть выполнены следующие условия\/$:$


\begin{itemize}


\item[1)] ``$t\to A(t)\text{''}\in W^{1,1}(0,T;\cal L(D;H));$


\item[2)] ``$t\to a(t,\cdot,\cdot\cdot)\text{''}\in W^{3,1}(0,T;\cal B(V));$


\item[3)] ``$t\to B(t)\text{''}\in W^{2,1}(0,T;\cal L(V;H));$


\item[4)] $f\in W^{2,1}(0,T;H);$


\item[5)] $f(0)\in V$.


\end{itemize}


Тогда $u\in W^{3,\infty}(0,T;H)\cap W^{2,\infty}(0,T;V)\cap W^{1,1}(0,T;D)$.


\end{thm}





\begin{thm}


Пусть выполнены следующие условия\/$:$


\begin{itemize}


\item[1)] ``$t\to A(t)\text{''} \in W^{2,1}(0,T;\cal L(D;H));$


\item[2)] ``$t\to a(t,\cdot,\cdot\cdot)\text{''}\in W^{4,1}(0,T;\cal B(V));$


\item[3)] ``$t\to B(t)\text{''}\in W^{3,1}(0,T;\cal L(V;H));$


\item[4)] $f\in W^{3,1}(0,T;H);$


\item[5)] $f(0)\in D$, $f'(0)\in V$.


\end{itemize}


Тогда $u\in W^{4,\infty}(0,T;H)\cap W^{3,\infty}(0,T;V)\cap W^{2,1}(0,T;D)$.


\end{thm}





Техника доказательства теорем 2, 3 близка используемой в


([5], с.\,30--32; [9], с.\,216--220).





\section{Схема проекционно-разностного метода}





Зададим последовательность $\{V_n\}_{n=1}^{\infty}$


конечномерных подпространств пространства $V$.


Примем обозначения: $A_n(t)$, $t\in[0,T]$, --- действующий в $V_n$ оператор,


определенный из тождества $(A_n(t)v_n,w_n)_H=a(t,v_n,w_n)$,


$v_n,w_n\in V_n$; $P_n$ и $Q_n$ --- ортопроекторы на $V_n$ соответственно


в $H$ и в $V$; $E$ --- тождественный оператор в $V$; $R_n=E-Q_n$.


Будем считать, что последовательность подпространств


$\{V_n\}_{n=1}^{\infty}$ предельно плотна в $V$,


т.\,е. для любого $v\in V$ $\|R_nv\|_V\to0$ при $n\to\infty$. Далее, обозначим


$\ov\omega_\tau=\{k\tau\}^{\infty}_{k=0}\cap[0,T]$


($\ov\omega_\tau$ --- сетка на отрезке $[0,T]$ с шагом $\tau$ и узлами


$k\tau$, $k=0,1,2,\dotsc$),


$\omega_\tau=\ov\omega_\tau\setminus\{0,\max\ov\omega_{\tau}\}$,


где $\tau$ --- любое число из $]0,T/2]$.





Зафиксируем некоторые


$\sigma_1, \sigma_2, \sigma_3, \sigma_4\in{\bold R}$,


$\sigma_1\ge\sigma_2$. Рассмотрим схему


проекционно-раз\-ност\-но\-го метода для


задачи (1.1), (1.2):


\begin{gather}


y_{\ov tt}+A_n(t)y^{(\sigma_1,\sigma_2)}(t)+


P_nB(t)y^{(\sigma_3,\sigma_4)}(t)=P_nf(t), t\in\omega_\tau,\\


y(t)\in V_n,\quad t\in\ov\omega_\tau,\ \; y(0)\text{ и } y(\tau)


\text{ заданы.}


\end{gather}


Здесь $n\in{\bold N}$, $\tau\in]0,T/2]$,


$y^{(\sigma_i,\sigma_{i+1})}=\sigma_i\wh y+(1-\sigma_i-\sigma_{i+1})y+


\sigma_{i+1}\check y$, $i=1$ или $i=3$.


Все рассматриваемые пары значений $(n,\tau)$ всюду в этой работе считаем


такими, что


\begin{equation}


\bigg(\frac{\sigma_1+\sigma_2}2-\frac{1+\varepsilon}4


\bigg)\tau^2A_n(t)+E\ge0, t\in\omega_\tau,


\end{equation}


где $\varepsilon>0$ --- константа, не зависящая от $n$, $\tau$ и $t$,


операторное неравенство понимается в смысле скалярного произведения


пространства $H$. Условия $\sigma_1\ge\sigma_2$ и (2.3) получаются, если


записать для главной части уравнения (2.1) (без члена


$P_nB(t)y^{(\sigma_3,\sigma_4)}(t)$) известные в теории операторно-разностных


схем условия устойчивости ([6], с.\,363; [7], с.\,235). Очевидно, если


$\sigma_1+\sigma_2>1/2$, то условие (2.3) выполнено при любых $n$ и


$\tau$ с константой $\varepsilon=2(\sigma_1+\sigma_2)-1$.





\begin{lem}


При любом выборе допустимой, т.\,е. удовлетворяющей


$(2.3)$, пары значений $(n,\tau)\in{\bold N}\times]0,T/2]$, где


$\tau$ достаточно мало $(\tau\le\tau_1$, где $\tau_1>0$ -- константа,


не зависящая от $n)$, и заданных начальных значений


$y(0), y(\tau)\in V_n$ задача $(2.1)$, $(2.2)$ имеет единственное


решение $y:\ov\omega_\tau\to V_n$.


\end{lem}





Справедливость этой леммы вытекает из того, что при любом выборе допустимой


пары $(n,\tau)$, где $\tau$ достаточно мало в указанном смысле, операторный


коэффициент $\tau^{-2}E+\sigma_1A_n(t)+\sigma_3P_nB(t)$ при $\wh y(t)$


в (2.1) положителен в смысле скалярного произведения пространства $H$ для


любого $t\in\omega_\tau$, в чем нетрудно убедиться, используя условия


$\sigma_1\ge \sigma_2$, (2.3) и неравенство (1.8).





Решение задачи (2.1), (2.2), соответствующее некоторой паре значений


$(n,\tau)$, удовлетворяющей требованиям леммы 1, и некоторым заданным


начальным значениям $y(0)$, $y(\tau)\in V_n$, будем трактовать как


приближенное решение задачи (1.1), (1.2), построенное проекционно-разностным


методом. Далее будем рассматривать следующие два варианта начальных условий


к схеме (2.1), (2.2):


\begin{equation}


y(0)=y(\tau)=0


\end{equation}


либо (при $f(0)\in V$)


\begin{equation}


y(0)=0,\quad y(\tau)=\frac{\tau^2}2Q_nf(0).


\end{equation}





Приведем нужную для дальнейшего оценку сеточных производных решения задачи


(2.1), (2.2) первого и второго порядков.





\begin{lem}


При любом выборе допустимой, т.\,е. удовлетворяющей


$(2.3)$, пары значений $(n,\tau)\in{\bold N}\times]0,T/2]$, где


$\tau$ достаточно мало $(\tau\le\tau_2$, $\tau_2\in\,]0,\tau_1]$ ---


константа, не зависящая от $n;$ $\tau_1$ --- константа, указанная в лемме


$1)$, для решения задачи $(2.1)$, $(2.2)$ с начальными условиями


$(2.4)$ либо при $f(0)\in V$ с начальными условиями $(2.5)$


выполнена оценка


\begin{equation}


\max\limits_{t\in\omega_\tau\cup\{0\}}\|y_t\|_V+


\max\limits_{t\in\omega_\tau}\|y_{\ov tt}\|_H\le M,


\end{equation}


где $M$ --- константа, не зависящая от $n$ и $\tau$.


\end{lem}





При доказательстве этой леммы используется техника сеточных энергетических


оценок ([6], сc.\,356--363, 370, 371;


[7], cс.\,225, 226, 233--236, 266, 267).





\section{Основные результаты}





Положим $\rho_n=\|R_n\|_{\cal L(D;V)}$, $n\in{\bold N}$.


Будем оценивать скорость сходимости схемы проекционно-разностного метода


(2.1), (2.2) для задачи (1.1), (1.2) в терминах величин $\rho_n$ и $\tau$. Из


принятых условий предельной плотности последовательности подпространств


$\{V_n\}_{n=1}^{\infty}$ в $V$


и компактности оператора $(A(0))^{-1}$ в $H$ следует, что $\rho_n\to0$ при


$n\to\infty$. Во многих случаях при конкретном выборе подпространств $V_n$


величину $\rho_n$ можно оценить сверху (по крайней мере асимптотически)


некоторой известной величиной, стремящейся к нулю при $n\to\infty$.


В конце статьи указана типичная асимптотическая оценка сверху


величины $\rho_n$ в приложении к первой начально-краевой задаче для


гиперболического уравнения второго порядка при дискретизации задачи


по пространству методов конечных элементов.





Оценки скорости сходимости схемы (2.1), (2.2) даны в теоремах 4--6. Теорема


4 относится к случаю естественной гладкости точного решения $u$ задачи


(1.1), (1.2), теоремы 5 и 6 --- к случаям повышенной гладкости точного решения.


Во всех этих теоремах имеется в виду, что $n\to\infty$, $\tau\to0$


таким образом, что пара значений $(n,\tau)$ остается допустимой


(удовлетворяющей (2.3)). В формулировке теоремы 6, где речь идет о схеме с


начальными условиями (2.5), условие $f(0)\in V$, необходимое для того, чтобы


второе начальное условие (2.5) имело смысл, опущено, т.\,к. oно является


следствием явно указанных условий этой теоремы.





\begin{thm}


Скорость сходимости схемы $(2.1)$, $(2.2)$


с начальными условиями $(2.4)$ характеризуется оценкой


\begin{equation}


\max\limits_{t\in\omega_{\tau}\cup\{0\},\,\delta\in[0,1]}[\|y_t-


u'(t+\delta\tau)\|_H+\|y(t)-u(t+\delta\tau)\|_V]


\le o(\rho_n^{1/2})+O(\tau^{1/2}).


\end{equation}


\end{thm}





\begin{thm}


Пусть $u\in W^{1,1}(0,T;D)$ и выполнено условие $1)$


теоремы $2$. Тогда скорость сходимости схемы $(2.1)$, $(2.2)$ с


начальными условиями $(2.4)$ характеризуется оценкой


\begin{equation}


\max\limits_{t\in\omega_\tau\cup\{0\},\,\delta\in[0,1]}[\|y_t-


u'(t+\delta\tau)\|_H+\|y(t)-u(t+\delta\tau)\|_V]


=O(\rho_n+\tau).


\end{equation}


\end{thm}





\begin{thm}


Пусть


$u\in W^{3,\infty}(0,T;H)\cap W^{2,\infty}(0,T;V)\cap W^{2,1}(0,T;D)$.


Пусть также выполнены условия $3)$ и $4)$ теоремы $2$ и


условие $1)$ теоремы $3$. Тогда скорость сходимости симметричной


$(\sigma_1=\sigma_2$, $\sigma_3=\sigma_4)$ схемы $(2.1)$, $(2.2)$


с начальными условиями $(2.5)$ характеризуется оценкой


\begin{equation}


\max\limits_{t\in\omega_\tau\cup\{0\}}\|y_t-u'(t+\tfrac{\tau}2)\|_H+


\max\limits_{t\in\ov\omega_{\tau}}\|y(t)-u(t)\|_V={O}(\rho_n+\tau^2).


\end{equation}


\end{thm}





\begin{pf*}{Доказательство теоремы 4}


Примем обозначения:


$Q_n(t)$ --- ортопроектор $V$ на $V_n$ в смысле скалярного произведения


$a(t,\cdot,\cdot\cdot)$, $R_n(t)=E-Q_n(t)$, $z(t)=y(t)-u(t)$


$(t\in\ov\omega_\tau)$,


$Z(t,\theta)=\|z_t\|_H^2+((\sigma_1+\sigma_2)/2-1/4)


\|Q_n(\theta)(\wh z(t)-z(t))\|_{(t)}^2+(1/4)


\|Q_n(\theta)(\wh z(t)+z(t))\|_{(t)}^2$


($t\in\omega_\tau\cup\{0\}$, $\theta=t$ или $\theta=t+\tau$),


$Z(t)=\|z_t\|_H^2+\|\wh z(t)\|_V^2+\|z(t)\|_V^2$


($t\in\omega_\tau\cup\{0\}$), $u^{(\sigma_i,\sigma_{i+1})}(t)=\sigma_i\wh u(t)


+(1-\sigma_i-\sigma_{i+1})u(t)+\sigma_{i+1}\check u(t)$


($t\in\omega_\tau$, $i=1$ или $i=3$). Далее полагаем $\tau\le\tau_2$, где


$\tau_2$ --- константа, указанная  в лемме 2. Через $c_i$, $i=1,2,\dotsc,20$,


ниже в этом доказательстве и в доказательствах теорем 5, 6 обозначим


различные положительные константы, не зависящие от рассматриваемых значений


$n$, $\tau$ и $t\in\omega_\tau$.





Заготовим для дальнейшего некоторые вспомогательные оценки. Пусть


$t$ --- любой узел сетки $\omega_\tau$, $v$ --- любой элемент пространства $V$.


На основе леммы Обэна--Нитше ([10], с.\,139) с помощью (1.3) и (1.6) получаем


\begin{equation}


\|R_n(t)v\|_H\le a_1a_2\rho_n\|R_nv\|_V.


\end{equation}


Отсюда, очевидно, следует


\begin{equation}


\|R_n(t)v\|_H\le a_1a_2\rho_n\|v\|_V.


\end{equation}


При имеющейся гладкости функции $u$ ($u\in W^{2,\infty}(0,T;H)


\cap W^{1,\infty}(0,T;V)$, см. теорему 1)


\begin{align}


\|u_{\ov tt}\|_H&\le c_1,\\


\max\limits_{\theta\in[t-\tau,t+\tau]}\|u(t)-u(\theta)\|_V&\le


c_2\tau,\\


\|u(t)-u^{(\sigma_i,\sigma_{i+1})}(t)\|_V&\le c_3\tau,\quad


i=1\text{ или }i=3,\\


\|u(\tau)\|_H&\le c_4\tau^2.


\end{align}


В силу (2.6), (3.6)


\begin{equation}


\|z_{\ov tt}\|_H\le c_5.


\end{equation}


Согласно (3.7)


\begin{align}


\|u_t\|_V&\le c_2, \\


%


\|u(\tau)\|_V&\le c_2\tau.


\end{align}


Из (3.7) также следует


\begin{equation}


\|u_{\overset{\,\circ}{t}}\|_V\le c_2.


\end{equation}


Из (3.5), (3.11) и (3.5), (3.13) соответственно следуют оценки


\begin{align}


\|R_n(t)u_t\|_H&\le c_6\rho_n,\\


\|R_n(t)u_{\overset{\,\circ}{t}}\|_H&\le c_6\rho_n.


\end{align}


Используя (3.15), получаем


\begin{align}


\|Q_n(t)z_{\overset{\,\circ}{t}}\|_H&\le\frac12


(\|z_t\|_H+\|z_{\ov t}\|_H)+c_6\rho_n,\\


\|Q_n(t)z_{\overset{\,\circ}{t}}\|_H&\le\max\limits_


{\theta\in\omega_{\tau}\cup\{0\}}\|z_t(\theta)\|_H+c_6\rho_n.


\end{align}


С помощью (1.5), (1.6), (2.6) и (3.13) выводим оценку


\begin{equation}


\|Q_n(t)z_{\overset{\,\circ}{t}}\|_V\le c_7.


\end{equation}





Перейдем непосредственно к выводу оценки (3.1). Заменим в уравнении (1.1)


обозначение $t$ на $t_1$ и почленно проинтегрируем это уравнение сначала по


$t_1\in[t_2,t_2+\tau]$, где $t_2$ ---


любое число из $[t-\tau,t]$, $t$ --- любой


узел сетки $\omega_{\tau}$, а затем по $t_2\in[t-\tau,t]$. Полученное


соотношение почленно умножим на $\tau^{-2}$ и вычтем из (2.1). Вновь полученное


равенство почленно умножим скалярно в $H$ на


$2\tau Q_n(t)z_{\overset{\,\circ}{t}}$ и, применяя приемы вывода основного


энергетического тождества в ([6], cс.\,356, 357), преобразуем к виду


\begin{equation}


Z(t,t)=Z(t-\tau,t)+\sum_{i=1}^7s_i+2\tau^2(\sigma_2-\sigma_1)


\|Q_n(t)z_{\overset{\,\circ}{t}}\|_{(t)}^2,


\end{equation}


где


\begin{align*}


s_1&=-2\tau(z_{\ov tt},R_n(t)u_{\overset{\,\circ}{t}})_H,\\


s_2&=2\bigg[\tau^{-1}\int_{t-\tau}^t


\int_{t_2}^{t_2+\tau}a(t_1,u(t_1),Q_n(t)z_{\overset{\,\circ}{t}})


dt_1dt_2-\tau a(t,u(t),Q_n(t)z_{\overset{\,\circ}{t}})\bigg],\\


s_3&=2\tau a(t,\,u(t)-u^{(\sigma_1,\sigma_2)}(t),\,Q_n(t)


z_{\overset{\,\circ}{t}}),\\


s_4&=2\bigg(\tau^{-1}\int_{t-\tau}^t\int_{t_2}^{t_2+\tau}B(t_1)u(t_1)


dt_1dt_2-\tau B(t)u(t),\,Q_n(t)z_{\overset{\,\circ}{t}}\bigg)_H,\\


s_5&=2\tau(B(t)(u(t)-u^{(\sigma_3,\sigma_4)}(t)),\,Q_n(t)


z_{\overset{\,\circ}{t}})_H,\\


s_6&=-2\tau(B(t)z^{(\sigma_3,\sigma_4)}(t),\,Q_n(t)


z_{\overset{\,\circ}{t}})_H,\\


s_7&=2\bigg(\tau f(t)-\tau^{-1}\int_{t-\tau}^t\int_{t_2}^{t_2+\tau}


f(t_1)dt_1dt_2,\,Q_n(t)z_{\overset{\,\circ}{t}}\bigg)_H.


\end{align*}


Применяя элементы техники вывода энергетического неравенства в ([6],


cс.\,362, 363) и лемму 2 из [11], учитывая при этом (1.5)--(1.7), получаем


\begin{equation}


Z(t-\tau,t)\le Z(t-\tau,t-\tau)+c_8\tau(\|z(t)\|^2_V+\|


\check z(t)\|_V^2).


\end{equation}


С помощью (3.10) и (3.4) оценим


\begin{equation}


s_1\le c_9\rho_n\int_{t-\tau}^{t+\tau}\|R_nu'(\theta)\|_Vd\theta.


\end{equation}


Представим $s_2$ в виде


$$


s_2=2\tau^{-1}\int_{t-\tau}^t\int_{t_2}^{t_2+\tau}\int_t^{t_1}


\frac d{d\theta}a(\theta,u(\theta),Q_n(t)


z_{\overset{\,\circ}{t}})d\theta\, dt_1dt_2.


$$


Учитывая (1.7), (1.6), тот факт, что


$u\in W^{1,\infty}(0,T;V)$, и (3.18), получим оценку


\begin{equation}


s_2\le c_{10}\tau^2.


\end{equation}


В силу (1.6), (3.8) с $i=1$ и (3.18) имеем


\begin{equation}


s_3\le c_{11}\tau^2.


\end{equation}


С учетом принятого в \S\,1 условия


``$t\to B(t)\text{''}\in W^{1,1}(0,T;\cal L(V;H))$,


используя его следствие ---


неравенство (1.8), тот факт, что


$u\in W^{2,\infty}(0,T;H)$, и неравенство (3.17), получаем


\begin{multline}


s_4 =2\tau^{-1}\bigg(\int_{t-\tau}^t\int_{t_2}^{t_2+\tau}\int_t^{t_1}


\frac d{d\theta}B(\theta)u(\theta)d\theta\, dt_1dt_2,\,Q_n(t)


z_{\overset{\,\circ}{t}}\bigg)_H\le\\


\le c_{12}\tau\bigg[\int_{t-\tau}^{t+\tau}\|B'(\theta)\|_


{\cal L(V;H)}d\theta+\tau\bigg](\max\limits_


{\theta\in\omega_{\tau}\cup\{0\}}\|z_t(\theta)\|_H+\rho_n).


\end{multline}


Из (1.8), (3.8) с $i=3$ и (3.17) следует


\begin{equation}


s_5\le c_{13}\tau^2(\max\limits_


{\theta\in\omega_{\tau}\cup\{0\}}\|z_t(\theta)\|_H+\rho_n).


\end{equation}


С помощью (1.8) и (3.16) выводим  неравенство


\begin{equation}


s_6\le c_{14}\tau(Z(t)+\check Z(t)+\rho_n\max\limits_


{\theta\in\ov\omega_{\tau}}\|z(\theta)\|_V).


\end{equation}


Учитывая принятое в \S\,1 условие


$f\in W^{1,1}(0,T;H)$ и применяя (3.17), получаем


\begin{multline}


s_7 =2\tau^{-1}\bigg(\int_{t-\tau}^t\int_{t_2}^{t_2+\tau}\int_{t_1}^t


f'(\theta)d\theta dt_1dt_2,\,Q_n(t)z_{\overset{\,\circ}{t}}\bigg)_H\le \\


\le\tau\int_{t-\tau}^{t+\tau}\|f'(\theta)\|_Hd\theta(\max\limits_


{\theta\in\omega_{\tau}\cup\{0\}}\|z_t(\theta)\|_H+c_6\rho_n).


\end{multline}





Оценим сверху правую часть (3.19), используя (3.20)--(3.27) и условие


$\sigma_1\ge\sigma_2$, в силу которого последнее слагаемое в правой части


(3.19) неположительно. Полученные неравенства, соответствующие


$t=\tau,2\tau,\dotsc,\eta$, где $\eta$ --- любой узел сетки


$\omega_{\tau}$, почленно сложим. Несколько увеличив правую часть вновь


полученного неравенства, имеем


\begin{multline}


Z(\eta,\eta) \le Z(0,0)+C_1\bigg[\rho_nr_n+\tau


+(\rho_n+\tau)\max\limits_{\theta\in\omega_{\tau}


\cup\{0\}}(Z(\theta))^{1/2}+ \\


+\tau\sum_{t=\tau}^{\eta}


(Z(t)+\check Z(t))\bigg],\quad \eta\in\omega_{\tau},


\end{multline}


где $r_n=\int\limits_0^T\|R_nu'(\theta)\|_Vd\theta$; через $C_i$,


$i=1,2,\dotsc,10$,


в (3.28) и ниже (в том числе в доказательствах теорем 5, 6)


обозначаются различные положительные константы, не зависящие от


$n$, $\tau$ и $\eta$. Применяя теорему Лебега ([12], с.\,302) с учетом


того, что $u\in W^{1,\infty}(0,T;V)$ и последовательность подпространств


$\{V_n\}_{n=1}^{\infty}$ предельно плотна в $V$, получаем


\begin{equation}


r_n\to0,\quad n\to\infty.


\end{equation}





Используя (2.3), (3.14), (1.5), (1.6) и тот факт, что


$u\in L^{\infty}(0,T;D)$ (см. теорему 1), выводим неравенство


\begin{equation}


Z(\eta)\le C_2(Z(\eta,\eta)+\rho_n^2), \eta\in\omega_{\tau}.


\end{equation}


Данное неравенство аналогично известным в теории разностных схем оценкам


снизу для составных норм ([6], с.\,370). В силу (3.9) и (3.12) при


рассматриваемых здесь начальных условиях (2.4)


\begin{align}


Z(0,0)&\le C_3\tau^2,\\


Z(0)&\le C_4\tau^2.


\end{align}


Из (3.28), (3.30), (3.31) следует


\begin{equation}


Z(\eta)\le C_5\bigg[\rho_nr_n+\rho_n^2+\tau+


(\rho_n+\tau)\max\limits_{\theta\in\omega_{\tau}\cup\{0\}}(Z(\theta)


)^{1/2}+\tau\sum_{t=\tau}^{\eta}(Z(t)+\check Z(t))\bigg],


\quad \eta\in\omega_{\tau}.


\end{equation}


Из (3.33), (3.32) с помощью леммы 4 из ([7], гл.\,III, \S\,1, с.\,171)


получаем, что


при достаточно малых $\tau$ (при $\tau\le\tau^{\ast}$, где


$\tau^{\ast}>0$ --- константа, не зависящая от $n$)


\begin{equation}


Z(\eta)\le C_6(\rho_nr_n+\rho_n^2+\tau),


\quad \eta\in\omega_{\tau}\cup\{0\}.


\end{equation}


При имеющейся гладкости функции $u$


\begin{equation}


\max\limits_{t\in\omega_{\tau}\cup\{0\},\,\delta\in[0,1]}\|u_t-u'(t+


\delta\tau)\|_H=O(\tau).


\end{equation}


Из (3.34) с учетом (3.35), (3.7) и (3.29) получаем (3.1).


\end{pf*}





\begin{pf*}{Доказательство теоремы 5}


Действуем по образцу доказательства теоремы 4, с той разницей, что вместо


(3.21)--(3.23) получаем и используем оценки


\begin{align}


s_1&\le c_9\rho_n^2\int_{t-\tau}^{t+\tau}\|u'(\theta)\|_Dd\theta, \\


s_2&\le c_{15}\tau\int_{t-\tau}^{t+\tau}(\|A'(\theta)\|_{\cal L(D;H)}+


\|u'(\theta)\|_D)d\theta(\max\limits_


{\theta\in\omega_{\tau}\cup\{0\}}\|z_t(\theta)\|_H+\rho_n),\notag\\


s_3&\le c_{16}\tau\int_{t-\tau}^{t+\tau}\|u'(\theta)\|_Dd\theta


(\max\limits_{\theta\in\omega_{\tau}\cup\{0\}}\|z_t(\theta)\|_H+


\rho_n)\notag


\end{align}


(в частности, (3.36) непосредственно следует из (3.21) по определению


$\rho_n$  при имеющемся условии $u\in W^{1,1}(0,T;D)$).


В результате вместо (3.34) выводим оценку


$Z(\eta)\le C_7(\rho_n^2+\tau^2)$,


$\eta\in\omega_{\tau}\cup\{0\}$, из которой с учетом (3.35) и (3.7)


получаем (3.2).


\end{pf*}





\begin{pf*}{Доказательство теоремы 6}


Вновь действуем в основном по образцу


доказательства теоремы 4. Для


$s_2$, $s_3$, $s_4$, $s_5$ и $s_7$ при условиях доказываемой теоремы удается


установить оценки


\begin{align}


s_2&\le c_{17}\tau^2\bigg[\int_{t-\tau}^{t+\tau}


(\|A''(\theta)\|_{\cal L(D;H)}+


\|u''(\theta)\|_D)d\theta+\tau\bigg](\max\limits_


{\theta\in\omega_{\tau}\cup\{0\}}\|z_t(\theta)\|_H+\rho_n),\\


s_3&\le c_{18}\tau^2\int_{t-\tau}^{t+\tau}


\|u''(\theta)\|_Dd\theta(\max\limits_


{\theta\in\omega_{\tau}\cup\{0\}}\|z_t(\theta)\|_H+\rho_n),\\


s_4&\le c_{19}\tau^2\bigg[\int_{t-\tau}^{t+\tau}


\|B''(\theta)\|_{\cal L(V;H)}d\theta+\tau\bigg](\max\limits_


{\theta\in\omega_{\tau}\cup\{0\}}\|z_t(\theta)\|_H+\rho_n),\\


s_5&\le c_{20}\tau^3(\max\limits_


{\theta\in\omega_{\tau}\cup\{0\}}\|z_t(\theta)\|_H+\rho_n),\\


s_7&\le \frac{\tau^2}3\int_{t-\tau}^{t+\tau}


\|f''(\theta)\|_Hd\theta(\max\limits_


{\theta\in\omega_{\tau}\cup\{0\}}\|z_t(\theta)\|_H+c_6\rho_n).


\end{align}


Для $Z(0,0)$, $Z(0)$ при рассматриваемых здесь начальных условиях (2.5)


получаем оценки


\begin{equation}


Z(0,0)\le C_8(\rho_n^2+\tau^4),


\quad Z(0)\le C_9(\rho_n^2+\tau^4),


\end{equation}


при этом используем формулу Маклорена для $u(\tau)$, условие


$u\in W^{3,\infty}(0,T;H)\cap W^{2,\infty}(0,T;V)$


и неравенство $\|R_nv\|_H\le\rho_n\|v\|_V$, $v\in V$


(следствие леммы Обэна--Нитше), где полагаем $v=f(0)$. Применяя (3.37)--(3.42)


вместо (3.22)--(3.25), (3.27), (3.31), (3.32), а также (3.36) вместо (3.21),


выводим вместо (3.34) оценку


\begin{equation}


Z(\eta)\le C_{10}(\rho_n^2+\tau^4),\quad


\eta\in\omega_{\tau}\cup\{0\}.


\end{equation}


В силу условия $u\in W^{3,\infty}(0,T;H)$


\begin{equation}


\max\limits_{t\in\omega_{\tau}\cup\{0\}}\|u_t-u'(t+\tfrac{\tau}2)\|_H


=O(\tau^2).


\end{equation}


Из (3.43) и (3.44) следует (3.3).


\end{pf*}





\begin{note}


Следствием (3.2), (3.3) является оценка


\begin{equation}


\max\limits_{t\in\ov\omega_{\tau}}\|y(t)-u(t)\|_V


=O(\rho_n+\tau^p),


\end{equation}


где $p=1$ при условиях теоремы 5, $p=2$ при условиях теоремы 6. Уже в случае


простейшего уравнения (1.1) с оператором $A$, не зависящим от $t$, и с


$B(t)\equiv0$ верно следующее: при любом выборе $\sigma_1$, $\sigma_2$ оценка


(3.45) с $p=1$ неулучшаема по порядку в классе точных решений, обладающих


гладкостью, требуемой в теореме 5, как оценка погрешности схемы (2.1), (2.2)


с начальными условиями (2.4); при любом выборе $\sigma\in{\bold R}$ оценка


(3.45) с $p=2$ неулучшаема по порядку в классе точных решений, обладающих


гладкостью, требуемой в теореме 6, как оценка погрешности симметричной схемы


(2.1), (2.2) с $\sigma_1=\sigma_2=\sigma$ и с начальными условиями (2.5)


(о выборе $\sigma_3$, $\sigma_4$ здесь ничего не говорится, т.\,к.


в указанном случае уравнение (2.1) не содержит члена


$P_nB(t)y^{(\sigma_3,\sigma_4)}(t)$). Тем самым и оценки (3.2), (3.3)


неулучшаемы по порядку в соответствующих классах вычислительных схем и


достаточно гладких точных решений.


\end{note}





\begin{note}


Как видно из проведенных доказательств, если


$f(0)\in V$, то теоремы 4 и 5 остаются в силе при замене начальных условий


(2.4) на (2.5).


\end{note}





Приведем типичную асимптотическую оценку сверху величины $\rho_n$ для случая,


когда (1.1), (1.2) --- первая начально-краевая задача для гиперболического


уравнения второго порядка в цилиндрической области $\Omega\times[0,T]$, где


$\Omega\subset{\bold R^m}$, $m\in{\bold N}$, ---


ограниченная область с гладкой границей, и дискретизация задачи по пространству


проводится методом конечных элементов. Будем использовать обозначения


пространств Соболева, принятые в ([10], с.\,22). В данном случае


$H=L^2(\Omega)$, $A(t)$ ($t\in[0,T]$) ---


эллиптические операторы второго порядка с гладкими коэффициентами, заданные на


$D=H^2(\Omega)\cap H_0^1(\Omega)$, так что $V=H_0^1(\Omega)$, нормы


$\|\cdot\|_D$ и $\|\cdot\|_{H^2(\Omega)}$


эквивалентны на $D$ в силу второго основного неравенства для эллиптических


операторов ([9], с.\,118), нормы $\|\cdot\|_V$ и $\|\cdot\|_{H^1(\Omega)}$


эквивалентны на $V$ в силу эллиптичности оператора $A(0)$. Для метода конечных


элементов характерно следующее аппроксимационное свойство подпространств


$V_n$ (напр., [10], с.\,135):


$$


\min\limits_{v_n\in V_n}\|v-v_n\|_{H^1(\Omega)}\le


Ch_n\|v\|_{H^2(\Omega)},


$$


где $v$ --- любая функция из


$H^2(\Omega)\cap H_0^1(\Omega)$, $C$ ---


константа, не зависящая от $v$ и $n$,


$h_n$ --- максимальный из диаметров конечных элементов, порождающих


подпространство $V_n$. Отсюда в силу сказанного об эквивалентности норм следует


$\rho_n={O}(h_n)$, $n\in{\bold N}$.


Совмещая эту оценку с основными результатами данной работы,


получаем оценки скорости сходимости схем проекционно-разностного метода


решения первой начально-краевой задачи для гиперболического уравнения второго


порядка с дискретизацией по пространству методом конечных элементов


в терминах величин $h_n$ и $\tau$.








\begin{thebibliography}{99}





\bibitem{1}


Злотник А.А. {\it Оценки скорости сходимости проекционно-сеточных методов


для гиперболических уравнений второго порядка} // Вычисл. процессы


и системы. -- M., 1991. -- \No\,8. -- С.\,116--167.


\bibitem{2}


Железовский С.Е., Букесова Н.Н. {\it Оценки погрешности проекционного


метода для абстрактного квазилинейного гиперболического уравнения}


// Изв. вузов. Математика. -- 1999. -- \No\,5. -- С.\,94--96.


\bibitem{3}


Смагин В.В. {\it Коэрцитивные оценки погрешностей проекционного и


проекционно-раз\-ност\-но\-го методов для параболических уравнений}


// Матем. сб. -- 1994. -- Т.\,185. -- \No\,11. -- С.\,79--94.


\bibitem{4}


Смагин В.В. {\it Оценки в сильных нормах погрешности


проекционно-разностного метода приближенного решения абстрактного


параболического уравнения} // Матем. заметки. -- 1997. -- Т.\,62.


-- Вып.\,6. -- С.\,898--909.


\bibitem{5}


Лионс Ж.-Л. {\it Некоторые методы решения нелинейных краевых задач.} --


М.: Мир, 1972. -- 588 с.


\bibitem{6}


Самарский А.А. {\it Теория разностных схем.} -- 3-е изд. -- М.: Наука,


1989. -- 616 с.


\bibitem{7}


Самарский А.А., Гулин А.В. {\it Устойчивость разностных схем.} -- М.:


Наука, 1973. -- 416 с.


\bibitem{8}


Крейн С.Г. {\it Линейные дифференциальные уравнения в банаховом


пространстве.} -- М.: Наука, 1967. -- 464 с.


\bibitem{9}


Ладыженская О.А. {\it Краевые задачи математической физики.} -- М.: Наука,


1973. -- 408 с.


\bibitem{10}


Сьярле Ф. {\it Метод конечных элементов для эллиптических задач.} -- М.:


Мир, 1980. -- 512 с.


\bibitem{11}


Смагин В.В. {\it Оценки погрешности полудискретных приближений


по Гал\"еркину для параболических уравнений с краевым условием типа Неймана}


// Изв. вузов. Математика. -- 1996. -- \No\,3. -- С.\,50--57.


\bibitem{12}


Колмогоров А.Н., Фомин С.В. {\it Элементы теории функций и функционального


анализа.} -- 5-е изд. -- М.: Наука, 1981. -- 544 с.





\end{thebibliography}





\vskip 2cm





\post{\it Саратовский государственный}{\it Поступила}


\post{\it социально-экономический университет}{20.09.2000}





\label{end}


\end{document}


